% Part: lambda-calculus
% Chapter: lambda-definability
% Section: introduction

\documentclass[../../../include/open-logic-section]{subfiles}

\begin{document}

\olfileid{lam}{ldf}{int}
\olsection{પરિચય}

પ્રથમ નજરે લૅમ્બડા કલન માત્ર વિધેયો અને તેમનો બીજા
પદો પર અનુપ્રયોગ દર્શાવતી અભિવ્યક્તિઓનું ખૂબ અમૂર્ત
કલન લાગે છે. તેની વાક્યરચનામાં આ વિધેયો કયા પ્રકારની
વસ્તુઓ પરનાં છે તેનો કોઈ નિર્દેશ નથી; ખાસ કરીને
પ્રાકૃતિક સંખ્યાઓનો ઉલ્લેખ નથી. તેની મૂળ ક્રિયાઓ—
અનુપ્રયોગ અને લૅમ્બડા અમૂર્તીકરણ—ફક્ત પ્રાકૃતિક
સંખ્યાઓ પરનાં નહીં, કોઈપણ વિધેયને લાગુ પડે છે.

તેમ છતાં થોડી સૂઝ વાપરીને લૅમ્બડા કલનમાં
અંકગણિતીય વિધેયો, એટલે પ્રાકૃતિક સંખ્યાઓ પરનાં
વિધેયો, વ્યાખ્યાયિત કરી શકાય છે. તે માટે દરેક
પ્રાકૃતિક સંખ્યા~$n \in \Nat$ને ખાસ
$\lambd$-પદ~$\num n$ આપીએ છીએ; તેને~$n$નો
\emph{ચર્ચ અંક} કહે છે. ચર્ચ અંકોનું નામ એલોન્ઝો
ચર્ચ પરથી પડ્યું છે.

\begin{defn}
  $n \in \Nat$ હોય તો તેને અનુરૂપ \emph{ચર્ચ અંક}
  $\num{n}$ એ~$n$નું નિરૂપણ છે:
  \[
    \num{n} \ident \lambd[fx][f^n(x)]
  \]
  અહીં $f^n(x)$ એટલે~$f$ને~$x$ પર $n$ વાર
  લાગુ કરવાનું પરિણામ. દાખલા તરીકે,
  $\num{0}$ એ $\lambd[fx][x]$ છે અને
  $\num{3}$ એ $\lambd[fx][f(f(f\,x))]$ છે.
\end{defn}
  
ચર્ચ અંક~$\num n$ એ લૅમ્બડા પદ તરીકે સંકેતિત
છે. તે બે આર્ગ્યુમેન્ટ~$f$ અને~$x$ સ્વીકારી
$f^n(x)$ પાછું આપતું વિધેય દર્શાવે છે. ચર્ચ
અંકો સ્પષ્ટપણે સામાન્ય સ્વરૂપમાં છે.

લૅમ્બડા કલનમાં પ્રાકૃતિક સંખ્યાઓનું નિરૂપણ ત્યારે જ
ઉપયોગી બને જો તેમની સાથે ગણતરી કરી શકીએ. ચર્ચ અંક
સાથે ગણતરી એટલે $\lambd$-પદ~$F$ને એ અંક પર
લાગુ કરીને સંયુક્ત પદ~$F\, \num n$નું સામાન્ય
સ્વરૂપમાં લઘુકરણ કરવું. આ હંમેશા સામાન્ય સ્વરૂપમાં
લઘુકૃત થાય અને તે સ્વરૂપ હંમેશા ચર્ચ અંક~$\num m$
હોય તો ગણતરીનું પરિણામ સંખ્યા~$m$ ગણી શકાય.
પછી~$F$ વિધેય $f\colon \Nat \to \Nat$ને
વ્યાખ્યાયિત કરે છે એમ કહી શકીએ; ચોક્કસ રીતે,
$f(n) = m$ ત્યારે અને માત્ર ત્યારે જ જ્યારે
$F\, \num n \red \num m$. ચર્ચ--રોસર ગુણધર્મને
કારણે સામાન્ય સ્વરૂપ અસ્તિત્વમાં હોય તો અનન્ય છે.
તેથી $F\, \num n \red \num m$ હોય તો
$F \, \num n$નું લઘુકરણ બીજા કોઈ સામાન્ય
સ્વરૂપમાં, ખાસ કરીને બીજા ચર્ચ અંકમાં, થઈ શકે નહીં.

ઊલટું, વિધેય $f\colon \Nat \to \Nat$ આપેલું હોય
તો આ રીતે વિધેય~$f$ને વ્યાખ્યાયિત કરતું પદ~$F$ છે કે
નહીં તે પૂછી શકાય. હોય તો~$F$ વિધેય~$f$ને
\emph{!!{lambda define}s} કરે છે અને~$f$
!!{lambda definable} છે એમ કહીએ. આ વિચારને
અનેક-સ્થાનીય અને આંશિક વિધેયો સુધી વિસ્તારી શકાય.

\begin{defn}
માનો કે $f\colon \Nat^k \pto \Nat$. લૅમ્બડા
પદ~$F$ એ~$f$ને \emph{!!{lambda define}s}
કરે છે એમ કહીએ, જો બધા
$n_0$, \dots,~$n_{k-1}$ માટે,
\[
F \, \num{n_0} \, \num{n_1} \dots \num{n_{k-1}} \red \num{f(n_0, n_1, \dots,
  n_{k-1})}
\]
જ્યારે $f(n_0, \dots, n_{k-1})$ વ્યાખ્યાયિત
હોય; અને બાકી કિસ્સામાં
$F \, \num{n_0} \, \num{n_1} \dots \num{n_{k-1}}$
ને કોઈ સામાન્ય સ્વરૂપ ન હોય.
\sourcecorrection{OLLAM-046}{સ્થિર અંગ્રેજી મૂળમાં
અહીં $f\colon \Nat^k \to \Nat$ લખાયું છે,
પરંતુ તરત જ $f$ અવ્યાખ્યાયિત હોય તેવો કિસ્સો
આપ્યો છે. આંશિક વિધેય માટે યોગ્ય સંકેત
$f\colon \Nat^k \pto \Nat$ અહીં વાપર્યો છે.}
\end{defn}

અચલ વિધેયો સરળ ઉદાહરણ છે. પદ
$C_k \ident \lambd[x][\num{k}]$ એ
$c_k\colon \Nat \to \Nat$ને !!{lambda define}s
કરે છે, જ્યાં $c_k(n) = k$. ખરેખર, દરેક~$n$
માટે $C_k \, \num n \ident
(\lambd[x][\num{k}])\num n \redone \num{k}$.
ઓળખ-વિધેય $\lambd[x][x]$ વડે
!!{lambda defined} છે. વધુ જટિલ વિધેયો વ્યાખ્યાયિત
કરવા માટે ઘણી વાર વધુ સૂઝ જોઈએ. તેથી દરેક
સંગણનીય વિધેય !!{lambda definable} છે એ કદાચ
આશ્ચર્યજનક લાગે. ઊલટું પણ સાચું છે: વિધેય
!!{lambda definable} હોય તો તે સંગણનીય છે.
\sourcecorrection{OLLAM-047}{સ્થિર અંગ્રેજી મૂળમાં
અચલ વિધેયને~$c_k$ નામ આપ્યા પછી શરત
$c(n)=k$ લખે છે. અહીં એ જ વિધેયનું નામ
$c_k(n)=k$ રાખ્યું છે.}

\end{document}
